
\documentclass[UTF8,11pt,fontset=windows]{ctexart}
\usepackage[paperwidth=240mm,paperheight=200mm,margin=14mm,top=12mm,bottom=12mm]{geometry}
\usepackage{amsmath,amssymb,mathtools,bm}
\usepackage{xcolor,needspace}
\pagestyle{empty}
\setlength{\parindent}{0pt}
\setlength{\parskip}{6pt}
\linespread{1.08}
\setlength{\emergencystretch}{3em}
\allowdisplaybreaks[2]
\xeCJKDeclareCharClass{CJK}{"2460 -> "24FF}
\xeCJKDeclareCharClass{CJK}{"2160 -> "2188}
\ctexset{section={format=\Large\bfseries,beforeskip=12pt,afterskip=7pt}}
\begin{document}
{\LARGE\bfseries 2007年考研数学二答案与解析\par}\vspace{6pt}

\begingroup\color{gray}我的解答，未与官方核对。\par\endgroup

\section*{答案速查}

\begingroup\color{gray}选择题：1.B　2.A　3.C　4.D　5.D　6.D　7.C　8.B　9.A　10.B。\par\endgroup

\begingroup\color{gray}填空题：11.$-1/6$；12.$1+\sqrt2$；13.$(-1)^n2^n n!/3^{n+1}$；14.$C_1\mathrm e^x+C_2\mathrm e^{3x}-2\mathrm e^{2x}$；15.$2v f_v-2u f_u$，其中 $u=y/x,v=x/y$；16.$1$。\par\endgroup

17. $f(x)=\ln(\sin x+\cos x)$。

18. $V(a)=\pi a^2/(\ln a)^2$，在 $a=\mathrm e$ 处取得最小值 $\pi\mathrm e^2$。

19. $y=\dfrac23x^{3/2}+\dfrac13$。

20. $z'(0)=0,z''(0)=1$。

21. 存在性证明见正文。

22. 积分为 $1/3+4\sqrt2\ln(1+\sqrt2)$。

23. $a=1$ 时公共解为 $t(1,0,-1)^{\mathrm T}$；$a=2$ 时公共解为 $(0,1,-1)^{\mathrm T}$。

24. $B=\begin{pmatrix}0&1&-1\\1&0&1\\-1&1&0\end{pmatrix}$，特征值为 $-2,1,1$。

\clearpage

\section*{一、选择题（1～10）}

\Needspace{5\baselineskip}\textbf{1.（2007.1）答案：B。}

两个无穷小等价的判据是它们之比趋于1。仅仅同阶、或比值趋于一个非零常数，还不能称为等价。以下逐项与 $\sqrt x$ 比较。

A 中 $1-\mathrm e^{\sqrt x}\sim-\sqrt x$，与 $\sqrt x$ 的比值为 $-1$，不等价。

B 中

$$
\ln\frac{1+x}{1-\sqrt x}
=\ln(1+x)-\ln(1-\sqrt x)
=\sqrt x+O(x),
$$

所以与 $\sqrt x$ 等价。

C 中 $\sqrt{1+\sqrt x}-1\sim\sqrt x/2$，比值为 $1/2$；D 中 $1-\cos\sqrt x\sim x/2$，是更高阶的无穷小。因此只有 B。

B的比值还可以直接按基本极限计算：

$$
\begin{aligned}
\frac{\ln(1+x)-\ln(1-\sqrt x)}{\sqrt x}
&=\frac{\ln(1+x)}x\sqrt x
-\frac{\ln(1-\sqrt x)}{-\sqrt x}\,(-1)\\
&\longrightarrow 1\cdot0+1=1.
\end{aligned}
$$

C可用有理化得到比值 $1/(\sqrt{1+\sqrt x}+1)\to1/2$；D的比值为 $[(1-\cos\sqrt x)/x]\sqrt x\to0$。这些计算分别确认了每个选项的系数和阶数。

\clearpage

\Needspace{5\baselineskip}\textbf{2.（2007.2）答案：A，$x=0$。}

第一类间断点要求左右极限都存在且有限，包括可去间断点和跳跃间断点。因此本题不能只找使分母为零的位置，还要分别检查两个方向的极限。

将函数写成两个因子的乘积：

$$
f(x)=\frac{\mathrm e^{1/x}+\mathrm e}{\mathrm e^{1/x}-\mathrm e}\cdot\frac{\tan x}{x}.
$$

在 $x\to0^+$ 时，第一个因子趋于1，第二个趋于1，所以右极限为1。

在 $x\to0^-$ 时，第一个因子趋于 $-1$，第二个仍趋于1，所以左极限为 $-1$。

因此0是第一类跳跃间断点。

在 $x=1$ 处指数分母为零，另外因子不为零，函数无界；在 $x=\pm\pi/2$ 处正切函数无界，也都是第二类间断点。故选 A。

其中，零点右侧的指数因子可先将分子分母同除以 $\mathrm e^{1/x}$：

$$
\frac{\mathrm e^{1/x}+\mathrm e}{\mathrm e^{1/x}-\mathrm e}
=\frac{1+\mathrm e^{1-1/x}}{1-\mathrm e^{1-1/x}}\longrightarrow1.
$$

左侧则 $\mathrm e^{1/x}\to0$，因子趋于 $\mathrm e/(-\mathrm e)=-1$。在1附近，$\mathrm e^{1/x}-\mathrm e\sim-\mathrm e(x-1)$，而分子趋于 $2\mathrm e$，所以是无穷间断；在 $\pm\pi/2$ 处，指数比值有限且非零，无法抵消正切的无界性。

\clearpage

\Needspace{5\baselineskip}\textbf{3.（2007.3）答案：C。}

\Needspace{6\baselineskip}\textbf{先区分几何面积与有向积分。}

上半圆的积分取正值，下半圆的积分取负值。半圆面积为 $\pi r^2/2$，题目给的是直径，因此两个半径分别为 $1/2$ 和1。

四段有向积分依次为

$$
\int_{-3}^{-2}f=\frac\pi8.
$$

$$
\int_{-2}^{0}f=-\frac\pi2.
$$

$$
\int_0^2f=\frac\pi2.
$$

$$
\int_2^3f=-\frac\pi8.
$$

\Needspace{6\baselineskip}\textbf{再从定义求 $F$ 的各个取值。}

由四段积分相加减，

$$
F(2)=\frac\pi2.
$$

$$
F(3)=\frac\pi2-\frac\pi8=\frac{3\pi}{8}.
$$

负方向的积分要注意上下限反向：

$$
F(-2)=\frac\pi2.
$$

$$
F(-3)=-\left(\frac\pi8-\frac\pi2\right)=\frac{3\pi}{8}.
$$

因此 $F(-3)=\frac34F(2)$，选 C。

具体而言，$F(-2)=\int_0^{-2}f=-\int_{-2}^0f=\pi/2$，不能因负半轴上的图形在 $x$ 轴下方，就误把 $F(-2)$ 也记为负数。四个函数值已求出，逐项代入选项即可确认只有C成立。

\clearpage

\Needspace{5\baselineskip}\textbf{4.（2007.4）答案：D。}

\begingroup\color{gray}错题：2007-4 对称差商存在推不出可导\par\endgroup

这里“极限存在”指有限的两侧极限。题目要求选错误命题，下面前三项都能由连续性或导数定义推出。

A 中若 $f(x)/x$ 有有限极限，则 $f(x)\to0$，再由连续性得 $f(0)=0$。

B 中有限极限迫使分子 $f(x)+f(-x)\to0$，由连续性得 $2f(0)=0$，所以也正确。

C 中由 A 知 $f(0)=0$，于是原商正是导数差商，导数存在。

D 可用反例 $f(x)=|x|$ 排除：分子 $f(x)-f(-x)$ 恒为零，极限存在，但零点左右导数分别为 $-1,1$，不可导。因此错误的是 D。

A的细节是：若 $f(x)/x\to L$，则 $f(x)=x[f(x)/x]\to0$。B同理，若 $[f(x)+f(-x)]/x\to L$，则分子趋于0，而连续性又使该分子趋于 $2f(0)$，所以 $f(0)=0$。

C在得到 $f(0)=0$ 后，便有

$$
f'(0)=\lim_{x\to0}\frac{f(x)-f(0)}x
=\lim_{x\to0}\frac{f(x)}x=L.
$$

D的差式只保留了函数的奇对称部分，偶对称部分可能不可导。反例 $|x|$ 恰为偶函数，故差式恒为零，但 $[|x|-|0|]/x$ 的左、右极限分别为 $-1,1$。

\clearpage

\Needspace{5\baselineskip}\textbf{5.（2007.5）答案：D，共3条。}

\begingroup\color{gray}错题：2007-5 斜渐近线卡在无穷减无穷\par\endgroup

\Needspace{6\baselineskip}\textbf{第一步：竖直渐近线——先看有限奇点。}

$\ln(1+\mathrm e^x)$ 在整个实轴连续，$1/x$ 只在 $x=0$ 无定义；$x\to0$ 时 $\frac1x$ 无界而另一项有限，所以 $x=0$ 是竖直渐近线。

\Needspace{6\baselineskip}\textbf{第二步：$x\to+\infty$ 按斜率、截距两步求。}

先求斜率

$$
k=\lim_{x\to+\infty}\frac{f(x)}x
=\lim_{x\to+\infty}\left(\frac1{x^2}+\frac{\ln(1+\mathrm e^x)}x\right)
=1,
$$

其中 $\dfrac{\ln(1+\mathrm e^x)}x$ 是 $\dfrac\infty\infty$，用洛必达得 $\dfrac{\mathrm e^x}{1+\mathrm e^x}\to1$。

再求截距，把 $x$ 写成 $\ln \mathrm e^x$，两项合成一个对数：

$$
\begin{aligned}
b&=\lim_{x\to+\infty}\left[\frac1x+\ln(1+\mathrm e^x)-x\right]\\
&=\lim_{x\to+\infty}\left[\frac1x+\ln\frac{1+\mathrm e^x}{\mathrm e^x}\right]\\
&=\lim_{x\to+\infty}\left[\frac1x+\ln(1+\mathrm e^{-x})\right]=0.
\end{aligned}
$$

$k=1,\ b=0$，所以这一端有斜渐近线 $y=x$。

\Needspace{6\baselineskip}\textbf{第三步：$x\to-\infty$ 同样两步走。}

$x\to-\infty$ 时 $\frac1x\to0$、$\mathrm e^x\to0$，故 $f(x)\to0$。于是

$$
k=\lim_{x\to-\infty}\frac{f(x)}x=\frac{0}{-\infty}=0,
$$

$$
b=\lim_{x\to-\infty}\left[f(x)-0\cdot x\right]=\lim_{x\to-\infty}f(x)=0.
$$

$k=0$ 说明这一端是水平渐近线 $y=0$——水平线是斜线的特例，同一套公式照样管用，不是"求不出斜渐近线"。

\Needspace{6\baselineskip}\textbf{结论：$x=0$、$y=x$、$y=0$，共 3 条，选 D。}

\Needspace{4\baselineskip}
{\bfseries 化简 $\infty-\infty$ 用的恒等式\par}
$\ln(1+\mathrm e^x)=\ln[\mathrm e^x(1+\mathrm e^{-x})]=x+\ln(1+\mathrm e^{-x})$，所以 $\ln(1+\mathrm e^x)-x=\ln(1+\mathrm e^{-x})\to0$。

两端必须\textbf{分别}求 $k,b$；$k$ 有限（含 $k=0$）才有渐近线。把 $k=0$ 误当成"求不出斜渐近线"，就会漏掉左端的 $y=0$，答案掉进 C。

\clearpage

\Needspace{5\baselineskip}\textbf{6.（2007.6）答案：D。}

由 $f''>0$ 可知 $f'$ 严格增加。令相邻差 $d_n=f(n+1)-f(n)$。

中值定理给出 $d_n=f'(\xi_n)$，其中 $n<\xi_n<n+1$。因此 $\xi_{n+1}>\xi_n$，从而 $d_{n+1}>d_n$。

若 $u_1<u_2$，则 $d_1>0$，所以

$$
u_n=u_1+\sum_{k=1}^{n-1}d_k
\ge u_1+(n-1)d_1\longrightarrow+\infty.
$$

故数列必发散，D 成立。

\Needspace{6\baselineskip}\textbf{两个反例分别核实另一种初始大小关系。}

取 $f(x)=\mathrm e^{-x}$，则 $f''=\mathrm e^{-x}>0$、$u_1=\mathrm e^{-1}>u_2=\mathrm e^{-2}$，但 $u_n\to0$，所以B不成立。

取 $f(x)=x^2-10x=(x-5)^2-25$（相当于把抛物线 $y=x^2$ 右移，加减常数不改变差分 $d_n$），则 $f''=2>0$、$u_1=-9>u_2=-16$，但 $u_n=n^2-10n\to+\infty$，所以A不成立。

由 $u_1<u_2$ 导出的线性下界 $u_1+(n-1)d_1$ 已经直接排除C。凸性使相邻增量越来越大，一旦初始增量为正，之后就至少按一个固定正步长增长。

\Needspace{4\baselineskip}
{\bfseries 画一条 $y=x^2$ 能直接判断吗\par}
能定方向，也能当排除工具，但不能用它"认证"某个选项。抛物线本身就\textbf{推翻}了 C（$x^2$ 满足 $u_1<u_2$ 却发散）；A、B 的条件是 $u_1>u_2$，抛物线不落在它们那一支，得另取 $\mathrm e^{-x}$、$x^2-10x$ 作反例。

不过举例正是选择题的标准打法：$x^2$ 排掉 C、$\mathrm e^{-x}$ 排掉 B、$x^2-10x$ 排掉 A，剩下 D 就是答案——\textbf{反例排干净 + 单选设定}已经是完整解法，不必再一般地证 D。真正不成立的只有"用一个例子\textbf{确认}某选项"。\textbf{反例可以证伪，特例不能证真。}

还要留意"真空成立"：$x^2$ 不满足 A、B 的前提 $u_1>u_2$，这两个选项在它上面只是条件为假时的自动成立——例子对它们\textbf{中立}，既不算支持也不算否定。所以 $x^2$ 一项只对 C（反例）与 D（相符）说话，靠它分辨不出答案。

\clearpage

\Needspace{5\baselineskip}\textbf{7.（2007.7）答案：C。}

\begingroup\color{gray}错题：2007-7 沿轴的信息推不出可微\par\endgroup

设 $r=\sqrt{x^2+y^2}$。原点可微的定义是存在常数 $A,B$，使

$$
f(x,y)-f(0,0)=Ax+By+o(r).
$$

C正好给出了 $A=B=0$ 的这一特殊情形，因此它是充分条件。

其余条件均可用

$$
f(x,y)=\begin{cases}xy/\sqrt{x^2+y^2},&(x,y)\ne(0,0),\\0,&(x,y)=(0,0)
\end{cases}
$$

说明不足。它在原点连续，沿坐标轴的相关偏导数恒为零，但沿 $y=x$ 有 $f(x,x)/\sqrt{2x^2}=1/2$，不趋于零，因此不可微。

\Needspace{6\baselineskip}\textbf{核实反例同时满足A、B、D。}

由 $2|xy|\le x^2+y^2$，反例函数满足 $|f(x,y)|\le r/2\to0$，所以在原点连续，满足A。

沿两条坐标轴，$f(x,0)=f(0,y)=0$，故B中的两个差商都为零。进一步，$f_x(x,0)=0$、$f_y(0,y)=0$，包括原点处的相应偏导，故D中的两个极限也为零。

若它可微，线性主部必由这两个偏导确定为零；但沿 $y=x\ne0$，余项除以 $r$ 恒为 $1/2$，与可微定义矛盾。D只控制坐标轴上的偏导，不能代替整个二维邻域中的控制。

\clearpage

\Needspace{5\baselineskip}\textbf{8.（2007.8）答案：B。}

\Needspace{6\baselineskip}\textbf{先把累次积分转换为区域描述。}

外层的 $x$ 从 $\pi/2$ 到 $\pi$，对每一个固定的 $x$，内层 $y$ 从正弦曲线上升到水平线 $y=1$。换序时应保持这一块区域不变，只改用水平截线描述。

原区域满足 $\pi/2\le x\le\pi$、$\sin x\le y\le1$。投影到 $y$ 轴得到 $0\le y\le1$。

在这个 $x$ 范围内，正弦单调减少，$y\ge\sin x$ 等价于

$$
x\ge\pi-\arcsin y.
$$

所以换序后为

$$
\int_0^1\mathrm dy\int_{\pi-\arcsin y}^{\pi} f(x,y)\,\mathrm dx,
$$

即 B。

反解边界时，$\arcsin y$ 取值在 $[0,\pi/2]$，不属于原题所用的第二象限部分。原边界对应的解应是 $x=\pi-\arcsin y$。

固定 $y\in(0,1)$，正弦在 $[\pi/2,\pi]$ 上递减，满足 $\sin x\le y$ 的点位于该交点右侧，因此右端始终为 $\pi$。例如 $y=0$ 时只剩 $x=\pi$；$y=1$ 时整段 $[\pi/2,\pi]$ 都被包含，积分限与这两种边界情形一致。

\clearpage

\Needspace{5\baselineskip}\textbf{9.（2007.9）答案：A。}

A 中三个向量直接相加，得到

$$
(\alpha_1-\alpha_2)+(\alpha_2-\alpha_3)+(\alpha_3-\alpha_1)=0.
$$

这给出系数不全为零的线性关系，因此该向量组线性相关。

\Needspace{6\baselineskip}\textbf{从坐标矩阵看其他选项为什么不同。}

原向量组三个向量线性无关，故新向量组是否相关，可以通过它们相对于原向量组的系数矩阵判断。对循环形式

$$
\beta_1=\alpha_1+c\alpha_2,\quad
\beta_2=\alpha_2+c\alpha_3,\quad
\beta_3=\alpha_3+c\alpha_1,
$$

系数矩阵及其行列式为

$$
C=\begin{pmatrix}1&0&c\\c&1&0\\0&c&1\end{pmatrix},
\qquad|C|=1+c^3.
$$

A对应 $c=-1$，行列式为0；B、C、D分别对应 $c=1,-2,2$，行列式为 $2,-7,9$，均非零。因此只有A相关。直接找到三个向量相加为零，是此题最快的判断方式。

\clearpage

\Needspace{5\baselineskip}\textbf{10.（2007.10）答案：B，合同但不相似。}

\Needspace{6\baselineskip}\textbf{先分清相似与合同所保持的量。}

相似变换保持全部特征值，而实对称矩阵的合同变换保持正、负、零惯性指数。特征值大小不同，并不妨碍两个正系数通过变量缩放互相转换。

矩阵 $A=3I-J$，其中 $J$ 为三阶全1矩阵，所以 $A$ 的特征值为 $3,3,0$。

矩阵 $B$ 的特征值为 $1,1,0$。两者都是实对称矩阵，且具有相同惯性指数：两个正平方项、一个零项，所以合同。

但它们的特征值不同，迹分别为6和2，因此不相似。

\Needspace{6\baselineskip}\textbf{具体求出 $A$ 的三个特征方向。}

令 $e=(1,1,1)^{\mathrm T}$，则 $Je=3e$，故 $Ae=0$。若 $v_1+v_2+v_3=0$，就有 $Jv=0$，从而 $Av=3v$。满足坐标和为零的向量组成二维子空间，因此 $A$ 的特征值确为 $3,3,0$。

正交对角化后，$A$ 的二次型成为 $3s_1^2+3s_2^2$。令 $t_1=\sqrt3s_1,t_2=\sqrt3s_2,t_3=s_3$，即可化为 $t_1^2+t_2^2$，与 $B$ 的二次型相同。这个缩放可逆，证明合同；迹不同则立即证明不相似。

\clearpage

\section*{二、填空题（11～16）}

\Needspace{5\baselineskip}\textbf{11.（2007.11）答案：$-\dfrac16$。}

分子中两个函数都与 $x$ 等价，直接用一阶等价无穷小相减会使主要项抵消。因此必须保留到三次项。

在 $x=0$ 处使用泰勒展开：

$$
\arctan x=x-\frac{x^3}{3}+O(x^5),
\qquad
\sin x=x-\frac{x^3}{6}+O(x^5).
$$

相减时，一次项消去，得到

$$
\arctan x-\sin x
=\left(-\frac13+\frac16\right)x^3+O(x^5)
=-\frac{x^3}{6}+O(x^5).
$$

所以

$$
\frac{\arctan x-\sin x}{x^3}
=-\frac16+O(x^2)\longrightarrow-\frac16.
$$

这里三次项的负号来自 $-1/3+1/6$，不能只比较两个函数各自的正负。

\clearpage

\Needspace{5\baselineskip}\textbf{12.（2007.12）答案：$1+\sqrt2$。}

\Needspace{6\baselineskip}\textbf{先求切线斜率，再利用垂直关系求法线斜率。}

对参数 $t$ 求导，得

$$
\frac{\mathrm dx}{\mathrm dt}
=-\sin t-2\cos t\sin t
=-\sin t(1+2\cos t),
$$

$$
\frac{\mathrm dy}{\mathrm dt}=\cos t.
$$

在 $t=\pi/4$ 处，两个参数导数都不为零，因而切线斜率存在且非零：

$$
k_T=\frac{\mathrm dy/\mathrm dt}{\mathrm dx/\mathrm dt}
=-\frac{\cos(\pi/4)}{\sin(\pi/4)[1+2\cos(\pi/4)]}
=-\frac1{1+\sqrt2}.
$$

法线与切线垂直，所以 $k_Nk_T=-1$，从而

$$
\boxed{k_N=-\frac1{k_T}=1+\sqrt2.}
$$

也可以直接用 $k_N=-(\mathrm dx/\mathrm dt)/(\mathrm dy/\mathrm dt)$；只有分母不为零时，才能将它写成有限斜率。

\clearpage

\Needspace{5\baselineskip}\textbf{13.（2007.13）答案：$\dfrac{(-1)^n2^n n!}{3^{n+1}}$。}

\begingroup\color{gray}错题：2007-13 求n阶导漏了阶乘与幂次\par\endgroup

将函数写成 $y=(2x+3)^{-1}$。连续求导的前几项为

$$
y'=-2(2x+3)^{-2},
\qquad
y''=2^2\cdot2!(2x+3)^{-3},
$$

$$
y'''=-2^3\cdot3!(2x+3)^{-4}.
$$

每求导一次，幂指数减1，同时乘上原来的幂指数与内函数导数2。因此猜测并可归纳证明

$$
y^{(n)}(x)=(-1)^n2^n n!(2x+3)^{-n-1}.
$$

事实上，再求一次导数便得到

$$
\begin{aligned}
y^{(n+1)}(x)
&=(-1)^n2^n n![-(n+1)]\cdot2(2x+3)^{-n-2}\\
&=(-1)^{n+1}2^{n+1}(n+1)!(2x+3)^{-n-2},
\end{aligned}
$$

与同一公式一致。最后代入 $x=0$：

$$
\boxed{y^{(n)}(0)=\frac{(-1)^n2^n n!}{3^{n+1}}.}
$$

\clearpage

\Needspace{5\baselineskip}\textbf{14.（2007.14）答案：$C_1\mathrm e^x+C_2\mathrm e^{3x}-2\mathrm e^{2x}$。}

\Needspace{6\baselineskip}\textbf{第一步：求对应齐次方程的通解。}

齐次方程 $y''-4y'+3y=0$ 的特征方程是

$$
r^2-4r+3=(r-1)(r-3)=0.
$$

它有两个不同实根1、3，所以齐次通解为

$$
y_h=C_1\mathrm e^x+C_2\mathrm e^{3x}.
$$

\Needspace{6\baselineskip}\textbf{第二步：求一个非齐次特解。}

右端为 $2\mathrm e^{2x}$，指数2不是特征根，因此设 $y_p=K\mathrm e^{2x}$，无需另乘 $x$。此时

$$
y_p'=2K\mathrm e^{2x},\qquad y_p''=4K\mathrm e^{2x}.
$$

代回原方程得到

$$
(4K-8K+3K)\mathrm e^{2x}=2\mathrm e^{2x},
$$

故 $-K=2$，即 $K=-2$。将齐次通解与特解相加：

$$
\boxed{y=C_1\mathrm e^x+C_2\mathrm e^{3x}-2\mathrm e^{2x}.}
$$

其中 $C_1,C_2$ 为任意常数。

\clearpage

\Needspace{5\baselineskip}\textbf{15.（2007.15）答案：$2v f_v-2u f_u$，其中 $u=y/x,\ v=x/y$。}

在 $x\ne0,y\ne0$ 的定义域中，令

$$
u=\frac yx,\qquad v=\frac xy,\qquad z=f(u,v).
$$

求 $x$ 偏导时固定 $y$，求 $y$ 偏导时固定 $x$。内函数的四个偏导分别为

$$
u_x=-\frac y{x^2},\quad v_x=\frac1y,
\qquad
u_y=\frac1x,\quad v_y=-\frac x{y^2}.
$$

由多元复合函数链式法则，

$$
z_x=f_u u_x+f_v v_x
=-\frac y{x^2}f_u+\frac1y f_v,
$$

$$
z_y=f_u u_y+f_v v_y
=\frac1x f_u-\frac x{y^2}f_v.
$$

分别乘以 $x,y$ 后再相减：

$$
\begin{aligned}
xz_x-yz_y
&=\left(-\frac yx f_u+\frac xy f_v\right)
 -\left(\frac yx f_u-\frac xy f_v\right)\\
&=-2\frac yx f_u+2\frac xy f_v\\
&=\boxed{2v f_v-2u f_u}.
\end{aligned}
$$

这里的 $f_u,f_v$ 都在 $(u,v)=(y/x,x/y)$ 处取值。虽然复合后的 $uv=1$，链式法则中仍须分别计算外函数对两个自变量的偏导。

\clearpage

\Needspace{5\baselineskip}\textbf{16.（2007.16）答案：$1$。}

矩阵 $A$ 的非零元只在第一条上副对角线上。按矩阵乘法计算，得到

$$
A^2=\begin{pmatrix}
0&0&1&0\\0&0&0&1\\0&0&0&0\\0&0&0&0
\end{pmatrix}.
$$

$$
A^3=\begin{pmatrix}
0&0&0&1\\0&0&0&0\\0&0&0&0\\0&0&0&0
\end{pmatrix}.
$$

例如 $(A^2)_{13}=A_{12}A_{23}=1$，$(A^3)_{14}=(A^2)_{13}A_{34}=1$；其他位置没有连续的非零乘积。

$A^3$ 只有第四列非零，且该列为 $(1,0,0,0)^{\mathrm T}$。其列空间是一维空间，因此

$$
\boxed{r(A^3)=1.}
$$

\clearpage

\section*{三、解答题（17～24）}

\Needspace{5\baselineskip}\textbf{17.（2007.17）答案：$f(x)=\ln(\sin x+\cos x)$。}

\Needspace{6\baselineskip}\textbf{第一步：由反函数积分确定初值。}

题中积分使用了 $f^{-1}(0)$，因此0属于反函数的定义域。记

$$
c=f^{-1}(0)\in[0,\pi/4],
\qquad f(c)=0.
$$

在原等式中取 $x=c$，左端上下限同为0，于是

$$
\int_0^c t\frac{\cos t-\sin t}{\sin t+\cos t}\,\mathrm dt=0.
$$

当 $0<t<\pi/4$ 时，$t>0$、$\cos t-\sin t>0$，且分母为正，所以被积函数严格为正。若 $c>0$，这个积分必为正，与等式矛盾。因此

$$
c=0,\qquad\boxed{f(0)=0}.
$$

这个初值来自原积分条件，不能仅因积分下限为0就直接写出。

\Needspace{6\baselineskip}\textbf{第二步：对变上限积分求导。}

$f$ 连续且严格单调，故反函数在其定义区间上连续，可以使用变上限积分求导公式。等式左端对 $x$ 的导数为

$$
\frac{\mathrm d}{\mathrm dx}\int_0^{f(x)}f^{-1}(t)\,\mathrm dt
=f^{-1}(f(x))f'(x)=xf'(x).
$$

等式右端的导数为

$$
x\frac{\cos x-\sin x}{\sin x+\cos x}.
$$

因而在 $0<x<\pi/4$ 上约去 $x$，得到

$$
f'(x)=\frac{\cos x-\sin x}{\sin x+\cos x}
=\frac{\mathrm d}{\mathrm dx}\ln(\sin x+\cos x).
$$

\Needspace{6\baselineskip}\textbf{第三步：积分并利用初值确定常数。}

由于区间上 $\sin x+\cos x>0$，有

$$
f(x)=\ln(\sin x+\cos x)+C.
$$

由连续性将等式延伸到 $x=0$，代入 $f(0)=0$ 得 $C=0$。所以

$$
\boxed{f(x)=\ln(\sin x+\cos x),\qquad0\le x\le\frac\pi4.}
$$

\Needspace{6\baselineskip}\textbf{第四步：回到原条件检验。}

所得函数在 $(0,\pi/4)$ 上导数为正，故在整个闭区间严格增加，且值域为 $[0,\ln\sqrt2]$，反函数及题中积分均有定义。原等式两边在 $x=0$ 处同为0，在区间内部导数相同，因此两边恒相等。端点 $x=\pi/4$ 处导数为0，不影响严格单调性。

\clearpage

\Needspace{5\baselineskip}\textbf{18.（2007.18）答案：$V(a)=\dfrac{\pi a^2}{(\ln a)^2}$，在 $a=\mathrm e$ 时取得最小值 $\pi\mathrm e^2$。}

\Needspace{6\baselineskip}\textbf{（Ⅰ）以圆盘截面法建立反常积分。}

固定 $x\ge0$，区域绕 $x$ 轴旋转后形成半径为 $y$ 的圆盘，截面积为 $\pi y^2$。而

$$
y^2=x\,a^{-x/a}=x\mathrm e^{-cx},
\qquad c=\frac{\ln a}{a}.
$$

由 $a>1$ 知 $c>0$。于是

$$
V(a)=\pi\lim_{R\to+\infty}\int_0^R x\mathrm e^{-cx}\,\mathrm dx.
$$

在有限区间上分部积分，取 $u=x$、$\mathrm dv=\mathrm e^{-cx}\mathrm dx$，得到

$$
\begin{aligned}
\int_0^R x\mathrm e^{-cx}\,\mathrm dx
&=\left[-\frac{x}{c}\mathrm e^{-cx}\right]_0^R
 +\frac1c\int_0^R\mathrm e^{-cx}\,\mathrm dx\\
&=-\frac Rc\mathrm e^{-cR}+\frac1{c^2}(1-\mathrm e^{-cR}).
\end{aligned}
$$

因 $c>0$，当 $R\to+\infty$ 时，$\mathrm e^{-cR}\to0$ 且 $R\mathrm e^{-cR}\to0$。后者可由洛必达法则求 $R/\mathrm e^{cR}$ 的极限得到。因此反常积分收敛，并且

$$
\boxed{V(a)=\frac\pi{c^2}=\frac{\pi a^2}{(\ln a)^2}.}
$$

\Needspace{6\baselineskip}\textbf{（Ⅱ）在参数范围 $a>1$ 内求全局最小值。}

由于 $h(a)=a/\ln a>0$，且 $V(a)=\pi[h(a)]^2$，使 $V$ 最小与使 $h$ 最小等价。

$$
h'(a)=\frac{\ln a-1}{(\ln a)^2}.
$$

分母恒正，故 $1<a<\mathrm e$ 时 $h'<0$，$a>\mathrm e$ 时 $h'>0$。函数先严格减少后严格增加，所以唯一的全局最小值点为 $a=\mathrm e$。代入得

$$
\boxed{V_{\min}=V(\mathrm e)=\pi\mathrm e^2.}
$$

\clearpage

\Needspace{5\baselineskip}\textbf{19.（2007.19）答案：$y=\dfrac23x^{3/2}+\dfrac13$，$x>0$。}

\Needspace{6\baselineskip}\textbf{第一步：利用方程不显含 $y$ 降阶。}

令 $p=y'$，则 $y''=\mathrm dp/\mathrm dx$，原方程化为

$$
\frac{\mathrm dp}{\mathrm dx}(x+p^2)=p,
\qquad p(1)=1.
$$

在初始点 $(x,p)=(1,1)$ 处，有 $p'(1)=1/2\ne0$。因此初值附近可以把 $x$ 看作 $p$ 的函数，取倒数得到

$$
\frac{\mathrm dx}{\mathrm dp}=\frac{x+p^2}{p},
$$

即关于 $x(p)$ 的一阶线性微分方程

$$
\frac{\mathrm dx}{\mathrm dp}-\frac1p x=p.
$$

\Needspace{6\baselineskip}\textbf{第二步：解线性方程并使用导数初值。}

初值附近 $p>0$，取积分因子 $1/p$。方程变为

$$
\frac1p\frac{\mathrm dx}{\mathrm dp}-\frac{x}{p^2}=1,
\qquad
\frac{\mathrm d}{\mathrm dp}\left(\frac xp\right)=1.
$$

积分得

$$
\frac xp=p+C,\qquad x=p^2+Cp.
$$

代入 $x=1,p=1$，得到 $C=0$，所以 $p^2=x$。由 $p(1)=1$ 选择正分支：

$$
y'=p=\sqrt x.
$$

负分支在初始点的导数为 $-1$，不满足题意；$p\equiv0$ 对应常值函数，也不满足导数初值。

\Needspace{6\baselineskip}\textbf{第三步：积分并使用函数值初值。}

$$
y=\int\sqrt x\,\mathrm dx
=\frac23x^{3/2}+C_1.
$$

代入 $y(1)=1$ 得 $C_1=1/3$，从而

$$
\boxed{y(x)=\frac23x^{3/2}+\frac13,\qquad x>0.}
$$

\Needspace{6\baselineskip}\textbf{第四步：直接代回检验。}

在 $x>0$ 上，$y'=\sqrt x$、$y''=1/(2\sqrt x)$，于是

$$
y''(x+y'^2)=\frac1{2\sqrt x}(x+x)=\sqrt x=y'.
$$

并且 $y(1)=y'(1)=1$。在包含1的区间上，$x>0$ 时右端 $p/(x+p^2)$ 光滑，初值解唯一；所得分支可延伸到整个 $(0,+\infty)$。零点处二阶导数无有限极限，不纳入这一经典解的区间。

\clearpage

\Needspace{5\baselineskip}\textbf{20.（2007.20）答案：$z'(0)=0,\ z''(0)=1$。}

\Needspace{6\baselineskip}\textbf{第一步：求隐函数在零点的函数值及前两阶导数。}

在 $y-x\mathrm e^{y-1}=1$ 中令 $x=0$，立即得到 $y(0)=1$。若记

$$
F(x,y)=y-x\mathrm e^{y-1}-1,
$$

则 $F_y(0,1)=1\ne0$，故在零点附近确实存在光滑隐函数 $y(x)$，且 $y>0$，后面的 $\ln y$ 有定义。

对原方程求一次导数：

$$
y'-\mathrm e^{y-1}-x\mathrm e^{y-1}y'=0,
$$

即

$$
(1-x\mathrm e^{y-1})y'=\mathrm e^{y-1}.
$$

代入 $x=0,y=1$，得 $y'(0)=1$。再对该等式求导：

$$
\begin{aligned}
&(1-x\mathrm e^{y-1})y''
 -(\mathrm e^{y-1}+x\mathrm e^{y-1}y')y'\\
&\hspace{25mm}=\mathrm e^{y-1}y'.
\end{aligned}
$$

移项可写成

$$
(1-x\mathrm e^{y-1})y''
=2\mathrm e^{y-1}y'+x\mathrm e^{y-1}(y')^2.
$$

因此 $y''(0)=2$。

\Needspace{6\baselineskip}\textbf{第二步：求内函数的前两阶导数。}

令 $u(x)=\ln y(x)-\sin x$，则

$$
u(0)=0,\qquad
u'=\frac{y'}y-\cos x,
$$

$$
u''=\frac{y''}y-\frac{(y')^2}{y^2}+\sin x.
$$

代入已经求出的 $y(0)=1,y'(0)=1,y''(0)=2$，得到

$$
u'(0)=1-1=0,\qquad u''(0)=2-1+0=1.
$$

\Needspace{6\baselineskip}\textbf{第三步：对外函数使用复合求导公式。}

由 $z=f(u)$，

$$
z'=f'(u)u',
\qquad
z''=f''(u)(u')^2+f'(u)u''.
$$

于是利用题设 $f'(0)=1$，

$$
\boxed{z'(0)=f'(0)u'(0)=0,}
$$

$$
\boxed{z''(0)=f''(0)\cdot0^2+1\cdot1=1.}
$$

虽然题目没有给出 $f''(0)$ 的数值，但该项乘有 $[u'(0)]^2=0$，因此不影响结果。

\clearpage

\Needspace{5\baselineskip}\textbf{21.（2007.21）证明：存在 $\xi\in(a,b)$，使 $f''(\xi)=g''(\xi)$。}

\Needspace{6\baselineskip}\textbf{第一步：把目标改写成差函数二阶导数为零。}

令 $h=f-g$。由题设，$h$ 在 $[a,b]$ 上连续，在 $(a,b)$ 内二阶可导，并且

$$
h(a)=h(b)=0.
$$

如果再能找到一个内点 $c$ 使 $h(c)=0$，就可以对三个零点连续使用两次罗尔定理。

\Needspace{6\baselineskip}\textbf{第二步：利用相等的最大值寻找内点零点。}

设两函数在 $(a,b)$ 内分别于 $u,v$ 处取得共同最大值 $M$，即

$$
f(u)=M,\qquad g(v)=M,\qquad u,v\in(a,b).
$$

因为 $g(u)\le M$、$f(v)\le M$，所以

$$
h(u)=M-g(u)\ge0,
\qquad
h(v)=f(v)-M\le0.
$$

若 $u=v$，则直接有 $h(u)=0$。若两点不同且上述任一不等式取等号，该点就是所求零点。否则两点处函数值严格异号，由介值定理，在 $u,v$ 之间存在零点 $c$。

无论哪一种情况，均得到

$$
c\in(a,b),\qquad h(c)=0.
$$

\Needspace{6\baselineskip}\textbf{第三步：在两个相邻区间分别用罗尔定理。}

在 $[a,c]$ 与 $[c,b]$ 上，$h$ 都连续，内部可导，且端点值相等。因此分别存在

$$
\alpha\in(a,c),\qquad\beta\in(c,b),
$$

使

$$
h'(\alpha)=0,\qquad h'(\beta)=0.
$$

\Needspace{6\baselineskip}\textbf{第四步：对一阶导函数再用一次罗尔定理。}

$h$ 在 $(a,b)$ 内二阶可导，故 $h'$ 在 $[\alpha,\beta]$ 上连续，在 $(\alpha,\beta)$ 内可导。又因两端值都为0，存在

$$
\xi\in(\alpha,\beta)\subset(a,b)
$$

使 $h''(\xi)=0$，即

$$
\boxed{f''(\xi)=g''(\xi).}
$$

这里用到的是函数在开区间内取得最大值，从而保证所找的第三个零点也是内点。

\clearpage

\Needspace{5\baselineskip}\textbf{22.（2007.22）答案：$\dfrac13+4\sqrt2\ln(1+\sqrt2)$。}

\begingroup\color{gray}错题：2007-22 分段函数配错区域又丢一层积分\par\endgroup

\Needspace{6\baselineskip}\textbf{第一步：按被积函数的两段定义拆分区域。}

记

$$
D_1=\{(x,y):|x|+|y|\le1\},
\qquad
D_2=\{(x,y):1<|x|+|y|\le2\}.
$$

所求积分为 $I=I_1+I_2$，其中

$$
I_1=\iint_{D_1}x^2\,\mathrm dx\,\mathrm dy,
\qquad
I_2=\iint_{D_2}\frac{\mathrm dx\,\mathrm dy}{\sqrt{x^2+y^2}}.
$$

两区域及被积函数均关于两条坐标轴对称，故只积第一象限再乘4。分界线面积为零，外环远离原点，无须另作反常积分处理。

\Needspace{6\baselineskip}\textbf{第二步：用直角坐标计算内层菱形上的积分。}

第一象限内 $D_1$ 是三角形 $0\le x\le1,\ 0\le y\le1-x$。于是

$$
\begin{aligned}
I_1
&=4\int_0^1\mathrm dx\int_0^{1-x}x^2\,\mathrm dy\\
&=4\int_0^1x^2(1-x)\,\mathrm dx\\
&=4\left(\frac13-\frac14\right)=\frac13.
\end{aligned}
$$

\Needspace{6\baselineskip}\textbf{第三步：用极坐标计算外层菱形环上的积分。}

第一象限令 $x=r\cos\theta,y=r\sin\theta$。边界 $x+y=1,2$ 分别变成

$$
r=\frac1{\cos\theta+\sin\theta},
\qquad
r=\frac2{\cos\theta+\sin\theta},
$$

而 $0\le\theta\le\pi/2$。极坐标面积元为 $r\,\mathrm dr\,\mathrm d\theta$，它与被积函数中的 $1/r$ 抵消。因此

$$
\begin{aligned}
I_2
&=4\int_0^{\pi/2}\mathrm d\theta
 \int_{1/(\cos\theta+\sin\theta)}^{2/(\cos\theta+\sin\theta)}\mathrm dr\\
&=4\int_0^{\pi/2}\frac{\mathrm d\theta}{\cos\theta+\sin\theta}.
\end{aligned}
$$

\Needspace{6\baselineskip}\textbf{第四步：完成三角积分并合并结果。}

利用 $\cos\theta+\sin\theta=\sqrt2\cos(\theta-\pi/4)$，令 $u=\theta-\pi/4$。其中正割积分可通过同乘 $\sec u+\tan u$ 凑微分：

$$
\mathrm d(\sec u+\tan u)
=\sec u(\tan u+\sec u)\,\mathrm du.
$$

因此

$$
\begin{aligned}
\int\sec u\,\mathrm du
&=\int\frac{\mathrm d(\sec u+\tan u)}{\sec u+\tan u}\\
&=\ln|\sec u+\tan u|+C.
\end{aligned}
$$

代入本题的积分限，得到

$$
\begin{aligned}
I_2
&=2\sqrt2\int_{-\pi/4}^{\pi/4}\sec u\,\mathrm du\\
&=2\sqrt2\left[\ln|\sec u+\tan u|\right]_{-\pi/4}^{\pi/4}\\
&=2\sqrt2\ln\frac{\sqrt2+1}{\sqrt2-1}.
\end{aligned}
$$

因为 $(\sqrt2-1)(\sqrt2+1)=1$，该比值等于 $(\sqrt2+1)^2$，所以

$$
I_2=4\sqrt2\ln(1+\sqrt2).
$$

最终

$$
\boxed{I=\frac13+4\sqrt2\ln(1+\sqrt2).}
$$

\clearpage

\Needspace{5\baselineskip}\textbf{23.（2007.23）答案：$a=1$ 时公共解为 $t(1,0,-1)^{\mathrm T}$；$a=2$ 时公共解为 $(0,1,-1)^{\mathrm T}$。}

\Needspace{6\baselineskip}\textbf{第一步：先分析齐次方程组何时允许公共解。}

记齐次方程组的系数矩阵为

$$
A=\begin{pmatrix}1&1&1\\1&2&a\\1&4&a^2\end{pmatrix}.
$$

用第二、三行分别减第一行，再用第三行减去第二行的3倍，得

$$
\begin{aligned}
A&\sim\begin{pmatrix}1&1&1\\0&1&a-1\\0&3&a^2-1\end{pmatrix}\\
&\sim\begin{pmatrix}1&1&1\\0&1&a-1\\0&0&(a-1)(a-2)\end{pmatrix}.
\end{aligned}
$$

这些行变换都不改变行列式，故

$$
\det A=(a-1)(a-2).
$$

若 $a\ne1,2$，则齐次方程组只有零解。但零解代入另一个方程要求 $0=a-1$，与 $a\ne1$ 矛盾。因此只需讨论 $a=1$ 与 $a=2$。

\Needspace{6\baselineskip}\textbf{第二步：讨论 $a=1$。}

阶梯形方程组化为

$$
x_1+x_2+x_3=0,\qquad x_2=0.
$$

所以 $x_1=-x_3$。另一个方程变为 $x_1+2x_2+x_3=0$，它已被满足，没有增加限制。于是所有公共解为

$$
\boxed{(x_1,x_2,x_3)^{\mathrm T}=t(1,0,-1)^{\mathrm T},\quad t\in\mathbb R.}
$$

\Needspace{6\baselineskip}\textbf{第三步：讨论 $a=2$。}

阶梯形方程组变为

$$
x_1+x_2+x_3=0,\qquad x_2+x_3=0,
$$

故 $x_1=0$、$x_3=-x_2$。代入另一个方程，得

$$
0+2x_2-x_2=2-1=1,
$$

从而 $x_2=1,x_3=-1$。此时唯一公共解为

$$
\boxed{(x_1,x_2,x_3)^{\mathrm T}=(0,1,-1)^{\mathrm T}.}
$$

代回原三行齐次方程分别得到 $0,0,0$，代回附加方程得到1，全部成立。以上已覆盖参数的所有情形。

\clearpage

\Needspace{5\baselineskip}\textbf{24.（2007.24）答案：$B=\begin{pmatrix}0&1&-1\\1&0&1\\-1&1&0\end{pmatrix}$，全部特征值为 $-2,1,1$。}

\Needspace{6\baselineskip}\textbf{（Ⅰ）通过矩阵多项式确定特征值。}

\begingroup\color{gray}错题：2007-24 矩阵多项式作用在特征向量上\par\endgroup

令 $p(t)=t^5-4t^3+1$，则 $B=p(A)$。若 $A\alpha=\lambda\alpha$，通过逐次相乘有 $A^k\alpha=\lambda^k\alpha$，因而

$$
B\alpha=p(A)\alpha=p(\lambda)\alpha.
$$

对题目给定的 $\alpha_1=(1,-1,1)^{\mathrm T}$，有 $A\alpha_1=\alpha_1$，所以

$$
B\alpha_1=(1-4+1)\alpha_1=-2\alpha_1.
$$

这就验证了 $\alpha_1$ 是 $B$ 属于特征值 $-2$ 的特征向量。

分别计算三个原特征值在多项式下的像：

$$
p(1)=-2,\qquad p(2)=32-32+1=1,
$$

$$
p(-2)=-32+32+1=1.
$$

$A$ 是实对称矩阵，可正交对角化，三个特征方向构成整个三维空间。因此 $B$ 的全部特征值为

$$
\boxed{-2,\ 1,\ 1.}
$$

\Needspace{6\baselineskip}\textbf{写出各特征值对应的全部特征向量。}

$A$ 的三个特征值不同，对应特征子空间两两正交。属于2与 $-2$ 的两个方向张成 $\alpha_1$ 的正交补；它们在 $B$ 作用下都乘以1，因而合并为 $B$ 的二维特征子空间。

属于 $-2$ 的全部特征向量为

$$
\boxed{k(1,-1,1)^{\mathrm T},\qquad k\ne0.}
$$

属于1的向量满足

$$
\alpha_1^{\mathrm T}x=x_1-x_2+x_3=0.
$$

令 $x_1=r,x_3=s$，便有 $x_2=r+s$，所以属于1的全部特征向量为

$$
\boxed{(r,r+s,s)^{\mathrm T},\qquad(r,s)\ne(0,0).}
$$

例如 $(1,1,0)^{\mathrm T}$ 与 $(0,1,1)^{\mathrm T}$ 构成这一特征子空间的一组基，但“全部特征向量”还包括它们的每个非零线性组合。

\Needspace{6\baselineskip}\textbf{（Ⅱ）设未知元素解方程。}

（Ⅰ）已经确定两件事：属于 $-2$ 的特征子空间是 $\alpha_1$ 张成的一维空间；属于 $1$ 的特征子空间是 $\alpha_1^{\perp}$。后者取两个线性无关的向量 $(1,1,0)^{\mathrm T}$ 与 $(1,0,-1)^{\mathrm T}$（都满足 $x_1-x_2+x_3=0$）。

$A$ 实对称时每个 $A^{k}$ 也对称——逐次转置即得 $(A^{k})^{\mathrm T}=(A^{\mathrm T})^{k}=A^{k}$——多项式各项相加，所以 $B=p(A)$ 仍实对称。故可设

$$
B=\begin{pmatrix}a&b&c\\b&d&e\\c&e&f\end{pmatrix},
$$

共六个未知量。把上面两件事写成方程：

$$
B\alpha_1=-2\alpha_1,
\qquad
B(1,1,0)^{\mathrm T}=(1,1,0)^{\mathrm T},
\qquad
B(1,0,-1)^{\mathrm T}=(1,0,-1)^{\mathrm T}.
$$

由 $B\alpha_1=-2\alpha_1$（$\alpha_1=(1,-1,1)^{\mathrm T}$）：

$$
a-b+c=-2,\qquad b-d+e=2,\qquad c-e+f=-2.
$$

由 $B(1,1,0)^{\mathrm T}=(1,1,0)^{\mathrm T}$：

$$
a+b=1,\qquad b+d=1,\qquad c+e=0.
$$

由 $B(1,0,-1)^{\mathrm T}=(1,0,-1)^{\mathrm T}$：

$$
a-c=1,\qquad b-e=0,\qquad c-f=-1.
$$

从 $a+b=1$ 与 $a-c=1$ 得 $b=1-a$、$c=a-1$，代入 $a-b+c=-2$：

$$
a-(1-a)+(a-1)=-2
\ \Longrightarrow\
3a-2=-2
\ \Longrightarrow\
a=0.
$$

于是 $b=1$、$c=-1$；再依次得 $e=1$、$d=0$、$f=0$。余下三条方程自动成立，可作核对：

$$
b-d+e=1-0+1=2,\qquad c-e+f=-1-1+0=-2,\qquad b-e=0.
$$

因此

$$
\boxed{B=\begin{pmatrix}0&1&-1\\1&0&1\\-1&1&0\end{pmatrix}.}
$$

\Needspace{4\baselineskip}
{\bfseries 另解：用正交谱分解直接写\par}
属于 $-2$ 与属于 $1$ 的两个特征子空间正交，把任意 $x$ 拆成 $x=Px+(I-P)x$，其中

$$
P=\frac{\alpha_1\alpha_1^{\mathrm T}}{\alpha_1^{\mathrm T}\alpha_1}
=\frac13\alpha_1\alpha_1^{\mathrm T}
$$

是到 $\alpha_1$ 方向的正交投影，则

$$
B=-2P+(I-P)=I-3P=I-\alpha_1\alpha_1^{\mathrm T}.
$$

把 $\alpha_1$ 与 $(1,1,0)^{\mathrm T}$ 分别乘回 $B$，得到 $-2$ 倍与原向量，可见特征值和特征向量都对上了；这与另解的表达式一致。
\end{document}
